\begin{lemma}[Claim 4.4: output of Algorithm (B)]
Retain the Case 1 data.  Let \(v_n\in V(T)\), and let the zero-based Lean
predecessor slot \(j_{\rm slot}\) satisfy the exact first-appearance identity
\[
 j_{\rm slot}.\mathrm{val}+1
   =\noderef{ETTFirstAppearanceIndex}(v_n)=j=m(v_n).
\]
Thus
\(\mathtt{ett.series.data.tree\ jSlot.succ.castSucc}\) is the exact stored
stage tree \(T_j\), while
\(\mathtt{ett.series.data.coloring\ jSlot.castSucc}\) is the exact reference
coloring \(\varphi_{j-1}\).  Choose a positive occurrence position of \(v_n\)
in this \(T_j\), certify that the vertex at that position is exactly \(v_n\),
and identify the closure seed by the equality
\[
\mathtt{seed}
   =\noderef{TreeSequencePrefix}(T_j,\mathtt{terminalPosition.val})
   =T_j(v_n).
\]
Let the Lean color-set parameter \(C\) be the manuscript set \(D_{j-1}\).
Write the support class in its manuscript order as
\(L_1=\{1_1<\cdots<1_{p+1}\}\), represented in Lean by
\(\mathtt{indices}:\operatorname{Fin}(p+1)\to\operatorname{Fin}(\mathtt{ett.rounds})\).
Retain the strict enumeration premise
\[
 a<c\quad\Longrightarrow\quad
 \mathtt{indices}(a).\mathrm{val}<\mathtt{indices}(c).\mathrm{val}.
\]
Require also that this enumeration has exactly the full support class as its
range: for every round slot \(r\),
\[
 r\in\noderef{ExtensionBaseSupportClass}
       (\mathtt{ett},v_n,\mathtt{hterminalSupport},v_n)
 \quad\Longleftrightarrow\quad
 \exists a\in\operatorname{Fin}(p+1),\ \mathtt{indices}(a)=r.
\]
Thus every active PE slot whose chosen supporting vertex is \(v_n\) occurs in
the displayed list, and every displayed slot is such an active PE slot.  This
is the complete class \(L_1\), including its first slot when the later selected
position is \(b=0\); it is not merely the processed prefix \(L_1^*\).
Also retain equation (4)'s distinctness assertion in its exact Lean form:
\[
 \mathtt{Fin.cons}(\gamma_0,\delta):
   \operatorname{Fin}(p+2)\longrightarrow\operatorname{Fin}(k)
 \quad\text{is injective}.
\]
Equivalently, the colors
\(\gamma_{1_1},\delta_{1_1},\ldots,\delta_{1_{p+1}}\) are pairwise
distinct; together with the chain identities, this is precisely the list
of distinct colors displayed in equation (4).  In particular,
\(\gamma_0=\gamma_{1_1}\) is the color named in (c3), and two nonadjacent
positions cannot collapse to the same trace slot.
Let \(T'_j\) be a \noderef{TashkinovClosure} of this exact ordered prefix under
\(\varrho_1\), let \(u\in V(T'_j)\setminus V(T_j)\), and let
\(\mathcal B\) be an
\noderef{ExtensionBaseAlgorithmBTrace} satisfying
\noderef{ExtensionBaseAlgorithmBInvariant}.  Write \(\varrho_1\) for
its initial coloring and \(\varrho_2\) for its output.  Assume Claim
4.3(b1), namely that \(\varrho_1\) is
\((T_j(v_n)-v_n,D_{j-1},\varphi_{j-1})\)-stable.  Assume the exact
Case 1 consequence of (14)--(15) saying that no color
used by Algorithm (B) is missing on \(T'_j(u)-u\), and that \(T'_j\) is a closure
of the displayed exact \(T_j(v_n)\) under \(\varrho_1\).  Finally retain the
least missing gamma at \(u\), the exact gamma and delta field equalities,
the exact chain identities, and the exact stored-pair identities.

Then:
\begin{enumerate}
\item[(c1)] \(\varrho_2\) is
 \((T_j(v_n)-v_n,D_{j-1},\varphi_{j-1})\)-stable;
\item[(c2)]
 \(\overline\varrho_2(v)=\overline\varrho_1(v)\) for every
 \(v\in V(T_j)\cup V(T'_j(u)-u)\), and
 \(\varrho_2(f)=\varrho_1(f)\) for every exact displayed edge identity
 \(f\) on \(T_j\) or on the prefix \(T'_j(u)\); and
\item[(c3)] \(\gamma_{1_1}\in\overline\varrho_2(u)\).
\end{enumerate}
In the Lean statement, \(V(T'_j(u)-u)\) is
\noderef{TreeVerticesBefore}, and an edge of \(T'_j(u)\) is expressed by
its exact edge index lying strictly before the unique vertex index of
\(u\).  Thus parallel edge identities are not collapsed.
\end{lemma}
\begin{proof}
Expand the three pieces of Algorithm (B) data first.  By
\noderef{ExtensionBaseCaseOneIndices}, the processed set is the image of
the positions \(q<b\).  By \noderef{ExtensionBaseAlgorithmBTrace}, if
\(r\) is the number of steps, then there are positions \(q_a<b\), exact
slots \(t_a=\mathtt{indices}(q_a)\), colors
\(\alpha_a,\beta_a\), and colorings
\[
 \varrho^0=\varrho_1,\varrho^1,\ldots,
 \varrho^r=\varrho_2                                      \tag{18a}
\]
such that the \(t_a\)'s have precisely that processed set as their image,
are strictly decreasing, and
\[
 \varrho^{a+1}=\noderef{TwoColorComponentSwap}
       (\varrho^a,\alpha_a,\beta_a,u).                      \tag{18b}
\]
The given strict increase of \(\mathtt{indices}\) makes it injective.
The strict decrease of the trace slots makes the trace slot map injective.
Consequently both sides of the trace's image equality have respectively
\(r\) and \(b\) elements, so \(r=b\).  Moreover, if \(a<c\), strict
decrease of the slots, their identities
\(t_a=\mathtt{indices}(q_a)\), and strict increase of
\(\mathtt{indices}\) give \(q_c<q_a\).  Thus the \(b\) distinct positions
below \(b\) occur in their unique decreasing order:
\[
 q_0=b-1,\ q_1=b-2,\ldots,\ q_{b-1}=0,\qquad
 t_a=\mathtt{indices}(q_a).                                \tag{18c}
\]
Here and below this display is omitted when \(b=0\).  The gamma and delta
field clauses of the trace, its slot-identity clause, and the hypotheses
on the corresponding ETT fields give, by injectivity of
\(\mathtt{some}\),
\[
                 \alpha_a=\gamma(q_a),\qquad
                 \beta_a=\delta(q_a).                      \tag{18d}
\]
This derives the exact slot and color identifications rather than merely
matching their numerical subscripts.

Now expand \noderef{ExtensionBaseAlgorithmBInvariant}.  Its first clause
says that every \(\varrho^a\), including the output, is proper in the
sense of \noderef{IsProperEdgeColoring}.  Before step \(a\), its second
clause says all of the following:
\[
 \begin{gathered}
 T_j\text{ is an ordered }\noderef{TreeSequenceInitialSegment}
       \text{ of }T_{t_a},\qquad
 \beta_a\in\overline{\varrho^a}(u),\\
 \noderef{TwoColorPathComponent}
       (G,\varrho^a,\alpha_a,\beta_a,u),                    \tag{18e}\\
 x\in V(T_{t_a})\text{ and }x\in
 \noderef{TwoColorComponent}(G,\varrho^a,
       \alpha_a,\beta_a,u)\Longrightarrow x=v_n,\\
 v_n\text{ in that component}\Longrightarrow
 v_n\text{ has at least two neighbors there}.
 \end{gathered}
\]
The last two lines say exactly that the path contains at most the one
vertex \(v_n\) from \(T_{t_a}\), and, when it contains \(v_n\), that
\(v_n\) is not an end.  Hence the manuscript's repeatedly used equation
is now justified in its exact form:
\[
 P_u(\alpha_a,\beta_a,\varrho^a)
 \text{ contains at most }v_n\text{ from }T_{t_a},
 \text{ and }v_n\text{ is not an end}.                     \tag{18}
\]

We will use the following elementary consequence of (18e).  The component
is the connected component through the vertex-node \(u\), by
\noderef{TwoColorComponent}; it contains an edge-node and its two colors
are distinct, by \noderef{TwoColorPathComponent}.  Since \(\beta_a\) is
missing at \(u\), no \(\beta_a\)-edge is incident with \(u\).  A path in
the connected component from \(u\) to the asserted edge-node must begin
with an incident edge-node.  By the definition of
\noderef{TwoColorIncidenceGraph}, that edge has color \(\alpha_a\).
Properness gives at most one incident edge of either fixed color, so this
is the unique neighbor of \(u\) in the two-color component.  Thus \(u\)
is an endpoint and has an \(\alpha_a\)-edge.

For completeness, the other endpoint is also a vertex-node.  Properness
gives every vertex-node at most one incident edge of each of the two
colors, hence degree at most two, as recorded by
\noderef{TwoColorIncidenceDegreeAtMostTwo}.  Every edge-node on the
component has exactly its two distinct endpoint vertices as neighbors:
its color is one of the selected colors and
\noderef{LooplessMultigraph} makes its two endpoints distinct.
In a finite connected graph of maximum degree two which has an edge and
an endpoint, following the unique unused neighbor from that endpoint must
terminate at one other endpoint; it cannot return to the first endpoint,
since the first has degree one, and it cannot continue forever.  Edge-nodes
have degree two, so the other endpoint is a vertex-node \(y\).  It has
exactly one incident path edge.  If that edge has color \(\alpha_a\), then
\(\beta_a\) is missing at \(y\); if it has color \(\beta_a\), then
\(\alpha_a\) is missing at \(y\).  Indeed, an incident edge of the other
color would be a second neighbor in the same component.

The same argument also describes every non-end vertex-node \(x\) in the
component.  Such an \(x\) has at least two neighbors and at most two.
The two incident edges cannot have the same color by properness, so one
has color \(\alpha_a\) and the other color \(\beta_a\).  The switch in
\noderef{TwoColorComponentSwap} exchanges these two values.  All incident
edges of other colors are fixed, and hence the set of colors present at
\(x\), and therefore \noderef{MissingColors}, is unchanged.  At the
endpoint \(u\), the same definition changes its unique
\(\alpha_a\)-edge to \(\beta_a\); because \(\beta_a\) was absent and
properness allowed only that one \(\alpha_a\)-edge, it makes
\(\alpha_a\) missing at \(u\).

We prove (c1).  Let \(x\) occur before \(v_n\) in the exact stage-\(j\)
tree order.  Then \(x\ne v_n\).  The ordered initial-segment clause in
(18e) places \(x\) in \(T_{t_a}\), so (18) says that \(x\) is not in the
switched component.  If an \(\alpha_a\)- or \(\beta_a\)-edge incident
with \(x\) were in that component, its adjacency to \(x\) in
\noderef{TwoColorIncidenceGraph} would put \(x\) in it as well.  Thus
such edges are outside the component; edges of all other colors are fixed
directly by \noderef{TwoColorComponentSwap}.  Every step therefore fixes
every edge incident with the before-\(v_n\) vertex set.  Induction on \(a\)
in (18a) gives
\[
 \varrho_2(f)=\varrho_1(f)\text{ on all those incident edges},
 \qquad
 \overline\varrho_2(x)=\overline\varrho_1(x)               \tag{19}
\]
for every such \(x\).  Expanding
\noderef{StableColoringBeforeVertex}, its assumed instance for
\(\varrho_1\) says that the controlled incident edges agree with
\(\varphi_{j-1}\) and that the missing sets agree with those under
\(\varphi_{j-1}\).  Composing those two equalities with (19) verifies the
two defining clauses for \(\varrho_2\).  This proves (c1).

We next prove the vertex half of (c2), simultaneously at every trace
stage.  At stage zero the desired equality with \(\varrho_1\) is
tautological.  Suppose it holds through stage \(a\).  If
\(x\in V(T_j)\) and \(x\ne v_n\), (18) makes the component avoid \(x\),
so the preceding incident-edge argument fixes its missing set.  If
\(x=v_n\), either the component avoids \(x\), with the same conclusion,
or the last clause of (18e) says that \(x\) is not an end.  The non-end
calculation above then says that the switch exchanges one edge of each
path color and again fixes its missing set.  Thus the induction preserves
the missing sets on all of \(V(T_j)\).

Now take \(x\in\noderef{TreeVerticesBefore}(T'_j,u)\).  If \(x\) is not
on the current component, the switch fixes its missing set.  If it is on
the component but is not internal, then \(x\ne u\), so it must be the
other endpoint just identified.  Consequently some
\(c\in\{\alpha_a,\beta_a\}\) is missing at \(x\) under \(\varrho^a\).
The simultaneous induction hypothesis converts this to
\(c\in\overline{\varrho_1}(x)\).  On the other hand, the trace enumeration
puts \(t_a\) in \noderef{ExtensionBaseCaseOneIndices}; (18c)--(18d) and
the stored-pair hypothesis identify
\[
 \mathtt{ett.series.data.stored}(t_a+1)
       =\{\alpha_a,\beta_a\}.
\]
Thus \(c\) belongs both to the missing-color union on the prefix and to
the stored-color union over the processed slots, contradicting
\(\mathtt{hprefixAvoids}\).  Therefore every prefix vertex which lies on
the component is internal.  The internal-vertex calculation above shows
that its present and missing sets survive the switch.  This completes the
induction and proves
\[
 \overline\varrho_2(x)=\overline\varrho_1(x)
 \quad\text{for every }x\in V(T_j)\cup V(T'_j(u)-u).         \tag{20}
\]

We turn to exact edge identities.  Fix a displayed edge \(f\) of \(T_j\)
and a trace step \(a\).  By \noderef{TreeSequence}, \(f\) joins its new
displayed vertex to a distinct earlier displayed vertex.  The
\noderef{TreeSequenceInitialSegment} clause places both vertices in
\(T_{t_a}\).  If the edge-node of \(f\) belonged to the switched
component, adjacency would put both of these distinct vertices in that
component, contrary to (18), which permits at most \(v_n\).  Hence the
edge-node is outside the component whenever its color is a path color;
if its color is not a path color, the switch definition fixes it directly.
Induction over (18a) proves
\[
                    \varrho_2(f)=\varrho_1(f)               \tag{21}
\]
for every exact displayed edge of \(T_j\).

It remains to prove (21) for an exact edge \(f\) of \(T'_j\) whose edge
index is before the displayed occurrence of \(u\).  From
\noderef{TashkinovClosure}, discard the final closedness clause and retain
the \noderef{TashkinovAugmentationSequence} from the exact seed to
\(T'_j\).  We spell out the needed induction through this
reflexive-transitive sequence.  For the empty sequence every final edge
is a seed edge.  For a sequence ending in one
\noderef{TashkinovAugmentation} \(Q\to Q^+\), that definition says that
\(Q^+\) has length \(\operatorname{length}(Q)+1\), that \(Q\) is an
ordered \noderef{TreeSequenceInitialSegment} of \(Q^+\), and that the
new last exact edge is the only edge not inherited from \(Q\).  Apply the
induction hypothesis to every inherited edge and record this last edge
with the immediately preceding tree \(Q\).  This proves the following
exhaustive alternative for each exact edge identity of \(T'_j\): it is an
edge inherited from the seed, or there is a unique augmentation stage at
which it is the new last edge.  Uniqueness follows because every step adds
one edge, every earlier tree is an ordered initial segment of every later
one, and the edge map of a \noderef{TreeSequence} is injective.

In the seed alternative, the equality
\(\mathtt{hseed}\) and the defining edge map of
\noderef{TreeSequencePrefix} identify \(f\) with the exact edge of
\(T_j\) at the same cast index.  Equation (21) applies.

In the augmentation alternative, let \(Q\to Q^+\) be the unique stage.
The definition of \noderef{TashkinovAugmentation} supplies a color \(c\)
such that
\[
              \varrho_1(f)=c,\qquad
              c\in\noderef{TreeMissingColors}(\varrho_1,Q). \tag{22}
\]
By the definition of \noderef{TreeMissingColors}, choose a displayed
vertex \(w\in V(Q)\) with
\(c\in\overline{\varrho_1}(w)\).  Chaining the ordered initial-segment
clauses identifies \(Q\) with the prefix of \(T'_j\) immediately before
the edge \(f\) was appended.  If the final edge index of \(f\) is \(q\),
then \(\operatorname{length}(Q)=q\), while every vertex of \(Q\) has
index at most \(q\).  The hypothesis on \(f\) says \(q\) is strictly less
than the vertex index of \(u\).  Hence the index of \(w\) is also strictly
less than that of \(u\), so
\[
              w\in\noderef{TreeVerticesBefore}(T'_j,u).      \tag{23}
\]
This is the missing intermediate-prefix witness used implicitly by TAA
in the manuscript.

Fix any processed switch \(a\).  As above, the exact trace slot belongs
to the processed set and its stored set is
\(\{\alpha_a,\beta_a\}\).  Equations (22)--(23) put \(c\) in the left
side of \(\mathtt{hprefixAvoids}\), while either equality
\(c=\alpha_a\) or \(c=\beta_a\) would put it in the right side.  Their
intersection is empty, so
\[
                       c\ne\alpha_a,\qquad c\ne\beta_a.      \tag{24}
\]
Induct now over the switches.  At stage zero (22) gives the exact value
of \(f\).  If \(\varrho^a(f)=c\), then (24) and the definition of
\noderef{TwoColorComponentSwap} show directly that
\(\varrho^{a+1}(f)=c\), whether or not one first tests membership of its
edge-node in the component.  Therefore
\(\varrho_2(f)=c=\varrho_1(f)\).  Together with the seed alternative this
proves the edge half of (c2), with no identification of parallel edges.

Finally prove (c3).  If \(b=0\), then
\noderef{ExtensionBaseCaseOneIndices} is empty.  The enumeration clause
of \noderef{ExtensionBaseAlgorithmBTrace} forces \(r=0\), since a
positive trace would contribute its first slot to the image.  The trace's
initial-coloring clause therefore identifies its output with
\(\varrho_1\).  The unique element \(b\) of \(\operatorname{Fin}(1)\) is
zero, so the local hypothesis \(\mathtt{hselected}\) is exactly
\(\gamma_0\in\overline{\varrho_1}(u)\), and hence gives (c3) directly.

Suppose \(b>0\).  By (18c), the first processed position is \(b-1\).
The chain identity gives
\(\delta(b-1)=\gamma(b)\), which is missing at \(u\) initially by
\(\mathtt{hselected}\).  More generally, before the step at position
\(q\), (18e) says that \(\delta(q)\) is missing at \(u\).  The endpoint
calculation following (18e), together with
\noderef{TwoColorComponentSwap}, says that this step makes
\(\gamma(q)\) missing at \(u\).  If \(q>0\), the chain identity rewrites
this color as \(\delta(q-1)\), exactly the missing delta needed for the
next position in the decreasing sequence (18c).  Descending induction
through \(b-1,b-2,\ldots,0\) therefore leaves
\(\gamma(0)=\gamma_{1_1}\) missing at \(u\) in the output coloring.
This is (c3), and (c1), (c2), and (c3) together prove the claim.
\end{proof}
